% ===========================================================================
%  benchmarks.tex
%
%  Literature-grounded benchmark suite for the complex-sinh Taylor-jet PINN.
%  Self-contained \section; class-agnostic (amsmath + booktabs only; uses
%  \eqref / \ref, no cleveref).  \input this into the main paper and add the
%  keys from benchmarks.bib.
%
%  Style follows the JCP/SIAM template in /root/paper/main.tex:
%    * governing equation, domain, and exact solution each in a display block;
%    * derivations (source terms) in align;
%    * prose kept separate from formulas;
%    * one compact "Setting" + "Provenance" note per family;
%    * booktabs tables.
% ===========================================================================

\section{Benchmark problems}\label{sec:benchmarks}

We evaluate the complex-parameter hyperbolic-sine network (``complex sinh''),
paired with the complex-Waring Taylor-jet residual backend of
Section~\ref{sec:impl}, on a suite of oscillatory and high-order partial
differential equations. Every problem is a benchmark drawn from, or directly
comparable to, the physics-informed neural network (PINN) literature; the exact
manufactured solutions and parameters used in our runs are stated in full below,
together with the source of each benchmark.

The suite is organized into three groups. The first probes \emph{high-frequency
second-order} problems, where spectral bias is the dominant difficulty
(Sections~\ref{sec:bm:helmholtz}--\ref{sec:bm:chirp}). The second probes
\emph{high-order real} operators, where the derivative order is the dominant
difficulty (Sections~\ref{sec:bm:poly}--\ref{sec:bm:ch}). The third probes
\emph{genuinely complex-valued} fields, where the unknown is naturally
\(\mathbb{C}\)-valued (Sections~\ref{sec:bm:nls}--\ref{sec:bm:maxwell}).
Table~\ref{tab:suite} summarizes the suite.

\begin{table}[htbp]
\centering
\small
\caption{Benchmark suite at a glance. Coordinates of a space--time problem are
written \((x,t)\); \(\Delta\) is the spatial Laplacian; \(a,k,m\) are the swept
parameters (wavenumber, mode, or temporal frequency). ``Source'' lists the
benchmark from which the setting is taken or to which it is matched.}
\label{tab:suite}
\begin{tabular}{@{}l l c l l@{}}
\toprule
family & operator & order & exact solution & source \\
\midrule
Helmholtz (high \(k\)) & \(\Delta u + \kappa^2 u\) & 2 & \(\sin(a\pi x)\sin(a\pi y)\) & \cite{Wang2021Pathologies} \\
Helmholtz (variable \(\kappa\)) & \(\Delta u + \kappa^2(x)u\) & 2 & \(\sin(a\pi x)\sin(a\pi y)\) & \cite{Hao2024PINNacle} \\
radial chirp & \(-\Delta u + u\) & 2 & \(\sin\!\big(\tfrac{a\pi}{2}r^2\big)\) & \cite{Tancik2020Fourier,Liu2020MscaleDNN} \\
polyharmonic & \(\Delta^m u\) & \(2m\) & \(\sin(\pi x)\,[\sin(\pi y)]\) & \cite{Vahab2022Biharmonic} \\
plate / beam & \(\Delta^2 w\) & 4 & \(\sin(m\pi x)\sin(n\pi y)\) & \cite{Vahab2022Biharmonic} \\
linearized KdV & \(u_t + \delta\,u_{xxx}\) & 3 & \(\sin(k\pi x)\cos(k\pi t)\) & \cite{Raissi2019} \\
Cahn--Hilliard & \(u_t - M\Delta(u^3-u-\gamma\Delta u\,[+\kappa\Delta^2 u])\) & 4/6 & \(\sin(a\pi x)\cos(a\pi t)\) & \cite{Raissi2019,Hao2024PINNacle} \\
cubic NLS & \(\mathrm{i}u_t + \tfrac12 u_{xx} + |u|^2 u\) & 2 & \(\operatorname{sech}(x)\,\mathrm{e}^{\mathrm{i}kt}\) & \cite{Raissi2019} \\
Maxwell (lossy) & \(\Delta E + \kappa^2 E\), \(\kappa^2\in\mathbb{C}\) & 2 & \(\mathrm{e}^{\mathrm{i}a\pi(x+y)}\) & \cite{Jiang2024Maxwell} \\
\bottomrule
\end{tabular}
\end{table}

% ---------------------------------------------------------------------------
\subsection{Shared protocol}\label{sec:bm:protocol}

All architectures and training choices are held fixed across the suite so that
the comparison isolates the effect of the network's value field
(\(\mathbb{C}\) versus \(\mathbb{R}\)) rather than incidental tuning.

\paragraph{Architectures.}
The complex-sinh network is a multilayer perceptron with \texttt{complex128}
parameters and the entire activation \(\sinh\). It is compared against four real
``high-frequency'' baselines that are standard remedies for spectral bias:
a random Fourier-feature network~\cite{Tancik2020Fourier}, a sinusoidal
representation network (SIREN)~\cite{Sitzmann2020SIREN}, a multi-scale deep
network (MscaleDNN)~\cite{Liu2020MscaleDNN}, and plain real \(\tanh\) and
\(\sinh\) baselines. For genuinely complex-valued targets
(Sections~\ref{sec:bm:nls}--\ref{sec:bm:maxwell}) each real baseline carries the
field as a split real--imaginary pair (an \(\mathbb{R}\)-valued PINN, RVPINN),
as in the original complex-valued PINN formulation~\cite{Raissi2019}.

\paragraph{Parameter matching.}
A width-\(H\) complex network has \(\approx 2H^2\) real degrees of freedom per
layer. To make the contest an equal-parameter one, every single real baseline is
widened to \(H_{\mathbb{R}} = \lceil \sqrt{2}\,H \rceil\) (so it has at least the
real parameter count of the complex network); for the split-real pairs each
component uses width \(H\) so that the \emph{pair} matches the complex network.

\paragraph{Frequency matching.}
Frequency-aware architectures are initialized at the problem's characteristic
wavenumber, so each method receives its best starting point at every swept
frequency: the SIREN / complex-sinh first-layer scale \(\omega_0\) and the
Fourier-feature bandwidth \(\sigma\) are set proportional to the target spatial
frequency of each case (specified per family below). For the high-order families
the initialization scale is held \emph{at} the target frequency rather than
above it, because an order-\(m\) operator amplifies an initialization frequency
\(\omega\) like \(\omega^{m}\); an over-large \(\omega_0\) is amplified by
\(\omega_0^{m}\) and buries the learning signal.

\paragraph{Derivative backend.}
All jet-compatible architectures (complex sinh, real \(\sinh/\tanh\), SIREN,
Fourier, MscaleDNN) evaluate every residual derivative with the \emph{same}
complex-Waring Taylor-jet operator of Section~\ref{sec:impl}, so the derivative
machinery does not bias the comparison. Nonlinear fluxes are written purely from
single-monomial partials of the network: for example the Cahn--Hilliard term
\begin{equation}\label{eq:lap_u3}
  \Delta(u^3) \;=\; 3u^2\,u_{xx} + 6u\,(u_x)^2
\end{equation}
uses only \(u\), \(u_x\), \(u_{xx}\), each a single Taylor-jet evaluation.

\paragraph{Optimization and metric.}
Training uses Adam with a wall-clock budget shared across all methods, a
time-based cosine learning-rate schedule decaying from \(10^{-3}\) to a floor of
\(10^{-4}\), and a fixed dense collocation set. The reported accuracy is the
relative \(L^2\) error of \(\operatorname{Re} u\) (or of \(u\) for complex-valued
targets) against the exact solution on a held-out grid of \(8192\) points,
\begin{equation}\label{eq:relL2}
  \mathcal{E}_{L^2}
  \;=\;
  \frac{\big\|\,u_\theta - u_\star\,\big\|_{L^2(\Omega)}}
       {\big\|\,u_\star\,\big\|_{L^2(\Omega)}} ,
\end{equation}
averaged over independent random seeds. Boundary and initial data are imposed by
a penalty term with weight \(100\); for the polyharmonic and plate families the
simply-supported (Navier) conditions are enforced on the lower even derivatives
(see Sections~\ref{sec:bm:poly}--\ref{sec:bm:plate}).

% ===========================================================================
\subsection{High-wavenumber Helmholtz}\label{sec:bm:helmholtz}

The first benchmark is the constant-coefficient Helmholtz equation on the unit
square,
\begin{equation}\label{eq:helm_pde}
  \Delta u + \kappa^2 u = f
  \qquad \text{in } \Omega = (-1,1)^2,
  \qquad \kappa = a\pi,
\end{equation}
with homogeneous Dirichlet data \(u = 0\) on \(\partial\Omega\). The exact
solution is the separable eigenmode
\begin{equation}\label{eq:helm_sol}
  u_\star(x,y) = \sin(a\pi x)\,\sin(a\pi y),
\end{equation}
and the source follows in closed form,
\begin{align}\label{eq:helm_source}
  f
  &= \Delta u_\star + \kappa^2 u_\star \notag\\
  &= \big(-2(a\pi)^2 + (a\pi)^2\big)\,u_\star
   = -(a\pi)^2\,u_\star .
\end{align}
The wavenumber is swept over \(a \in \{2,4,6,8,10\}\), so the number of
half-periods across the domain grows while the operator order stays fixed at
two; this isolates the effect of frequency on each architecture.

\paragraph{Provenance and comparability.}
This is the Helmholtz benchmark of Wang, Teng and
Perdikaris~\cite{Wang2021Pathologies}, who use
\(\Delta u + k^2 u = q\) on \((-1,1)^2\) with
\(u = \sin(a_1\pi x)\sin(a_2\pi y)\) and the matching source
\(q = (k^2 - (a_1\pi)^2 - (a_2\pi)^2)\,u\); the same family is the
Poisson--Boltzmann/Helmholtz task in the PINNacle suite~\cite{Hao2024PINNacle}.
We take the isotropic case \(a_1 = a_2 = a\) with \(k = a\pi\), and additionally
report the anisotropic case \((a_1,a_2) = (1,4)\) of~\cite{Wang2021Pathologies}
verbatim (with \(k=1\), \(q = (k^2 - (a_1\pi)^2 - (a_2\pi)^2)u_\star\)) as a
one-to-one numerical comparison.

\paragraph{Setting.}
\(\Omega=(-1,1)^2\); Dirichlet \(u=0\); isotropic sweep \(a\in\{2,4,6,8,10\}\)
plus the anisotropic \((1,4)\) case; \(\omega_0 = \max(10,\,2\pi a)\),
\(\sigma = \max(2,\,\pi a)\); residual scale \((a\pi)^2\).

% ===========================================================================
\subsection{Variable-coefficient (scattering) Helmholtz}\label{sec:bm:helmvc}

To move from a single clean eigenmode to a heterogeneous medium, we let the
zeroth-order coefficient vary in space,
\begin{equation}\label{eq:helmvc_pde}
  \Delta u + \kappa^2(x)\,u = f
  \qquad \text{in } \Omega = (-1,1)^2,
\end{equation}
with a smooth lens-like contrast
\begin{equation}\label{eq:helmvc_kappa}
  \kappa^2(x,y) = (a\pi)^2\big(1 + C\,\sin(\pi x)\sin(\pi y)\big),
  \qquad C = 0.5 .
\end{equation}
The manufactured solution and Dirichlet data are again
\eqref{eq:helm_sol} with \(u=0\) on \(\partial\Omega\), and the source is
\begin{equation}\label{eq:helmvc_source}
  f = -2(a\pi)^2\,u_\star + \kappa^2(x)\,u_\star .
\end{equation}
The background wavenumber is swept over \(a \in \{2,4,6\}\).

\paragraph{Provenance and comparability.}
Variable-coefficient (scattering) Helmholtz problems in a heterogeneous medium
are the realistic electromagnetics/acoustics regime represented in the PINNacle
Poisson--Boltzmann/Helmholtz tasks~\cite{Hao2024PINNacle}; the \(\pm 50\%\)
smooth contrast here is a controlled, exactly-known-solution instance of that
regime. It also exercises the variable-coefficient path of the residual
assembly, in which \(\kappa^2(x)\) is precomputed once on the collocation set.

\paragraph{Setting.}
\(\Omega=(-1,1)^2\); Dirichlet \(u=0\); contrast \(C=0.5\);
sweep \(a\in\{2,4,6\}\); \(\omega_0,\sigma\) as in
Section~\ref{sec:bm:helmholtz}.

% ===========================================================================
\subsection{Non-separable radial chirp}\label{sec:bm:chirp}

The eigenmode benchmarks above are separable products of sines, which
sine-based real baselines can represent almost natively and which may therefore
mask any expressivity gap. To remove this confound we use a non-separable target
with a space-varying local frequency,
\begin{equation}\label{eq:chirp_pde}
  -\Delta u + u = f
  \qquad \text{in } \Omega = (-1,1)^2,
\end{equation}
\begin{equation}\label{eq:chirp_sol}
  u_\star(x,y) = \sin\!\big(\phi(x,y)\big),
  \qquad
  \phi(x,y) = \tfrac{a\pi}{2}\,(x^2 + y^2),
\end{equation}
a radial chirp whose instantaneous frequency grows with radius,
\begin{align}\label{eq:chirp_freq}
  |\nabla\phi| = a\pi\,r, \qquad r = \sqrt{x^2+y^2}.
\end{align}
The source is assembled from
\begin{align}\label{eq:chirp_source}
  \Delta u_\star
  &= -(a\pi)^2\,r^2 \sin\phi + 2a\pi\,\cos\phi, \notag\\
  f &= -\Delta u_\star + u_\star,
\end{align}
with Dirichlet data \(u = u_\star\) on \(\partial\Omega\) and the sweep
\(a \in \{2,4,6,8\}\).

\paragraph{Provenance and comparability.}
Space-varying-frequency (chirp) targets are the standard stress test for
high-frequency representation methods such as Fourier
features~\cite{Tancik2020Fourier} and multi-scale
networks~\cite{Liu2020MscaleDNN}, which are tuned for stationary spectra. This
benchmark is the fairness control of the suite: every architecture must
\emph{construct} the varying oscillation rather than read it off a fixed
frequency basis.

\paragraph{Setting.}
\(\Omega=(-1,1)^2\); Dirichlet \(u=u_\star\); sweep \(a\in\{2,4,6,8\}\);
\(\omega_0,\sigma\) as in Section~\ref{sec:bm:helmholtz}.

% ===========================================================================
\subsection{Polyharmonic order sweep}\label{sec:bm:poly}

To isolate the effect of the differential \emph{order} from that of frequency,
we fix the solution and frequency and vary only the order of a polyharmonic
operator. The eigenmode identity \(\Delta^m u = (-S)^m u\) holds for the
sinusoidal eigenfunctions below, giving
\begin{equation}\label{eq:poly_pde}
  \Delta^m u = (-S)^m\,u,
  \qquad
  \begin{cases}
    u_\star = \sin(\pi x)\sin(\pi y), & S = 2\pi^2 \ (\text{2D}),\\[2pt]
    u_\star = \sin(\pi x),            & S = \pi^2 \ (\text{1D}).
  \end{cases}
\end{equation}
Simply-supported (Navier) boundary conditions are imposed exactly through the
lower even derivatives,
\begin{equation}\label{eq:poly_bc}
  \Delta^j u = 0 \quad \text{on } \partial\Omega,
  \qquad j = 0,1,\dots,m-1 .
\end{equation}
The order is swept over \(2m \in \{2,4,6\}\) in two dimensions and
\(2m \in \{2,4,6,8,10\}\) in one dimension. In one dimension the operator is a
single monomial partial \(\partial_x^{2m}\), which is cheap enough to reach order
ten; in two dimensions \(\Delta^m\) expands into \(m+1\) high-order terms.

\paragraph{Provenance and comparability.}
Biharmonic and polyharmonic operators are the canonical high-order PINN
testbed~\cite{Vahab2022Biharmonic}; the biharmonic (\(m=2\)) and triharmonic
(\(m=3\)) cases coincide with the polyharmonic family used in
randomized-smoothing PINNs. Holding \(u_\star\), the frequency and the domain
fixed while varying only \(m\) is our isolation design for the order axis.

\paragraph{Setting.}
\(\Omega=(-1,1)^d\), \(d\in\{1,2\}\); Navier BC \eqref{eq:poly_bc};
1D \(\omega_0=\sigma=\pi\), 2D \(\omega_0=2\pi\), \(\sigma=\pi\);
residual scale \(S^m\).

% ===========================================================================
\subsection{Kirchhoff plate and Euler--Bernoulli beam}\label{sec:bm:plate}

The fourth-order vibration eigenproblems of structural mechanics give two
oscillatory biharmonic benchmarks. For the Kirchhoff plate,
\begin{equation}\label{eq:plate_pde}
  \Delta^2 w = S^2\,w
  \qquad \text{in } \Omega = (-1,1)^2,
  \qquad S = (m^2 + n^2)\pi^2,
\end{equation}
with simply-supported edges \(w = 0\) and \(\Delta w = 0\) on \(\partial\Omega\)
and the mode shape
\begin{equation}\label{eq:plate_sol}
  w_\star(x,y) = \sin(m\pi x)\,\sin(n\pi y).
\end{equation}
For the Euler--Bernoulli beam in one dimension,
\begin{equation}\label{eq:beam_pde}
  w'''' = (m\pi)^4\,w
  \qquad \text{in } (-1,1),
  \qquad w_\star = \sin(m\pi x),
\end{equation}
with simply-supported ends \(w = 0\) and \(w'' = 0\). We sweep the mode number,
using the isotropic plate modes \((m,m)\) with \(m\in\{1,2,3\}\), the anisotropic
mixed modes \((m,m+1)\) with \(m\in\{2,3,4\}\), and the beam modes
\(m\in\{1,2,3\}\). The order is fixed at four; the sweep raises the oscillation
only.

\paragraph{Provenance and comparability.}
The biharmonic plate and beam are standard PINN structural benchmarks: Vahab et
al.~\cite{Vahab2022Biharmonic} solve the simply-supported Kirchhoff plate
(\(D\,\Delta^2 w = q\)), and the simply-supported Euler--Bernoulli beam with
mode shape \(\sin(m\pi x)\) is the textbook 1D specialization. We use the
free-vibration eigenmode form \(\Delta^2 w = S^2 w\) with the corresponding
manufactured right-hand side, so the exact solution is known in closed form.

\paragraph{Setting.}
Plate \(\Omega=(-1,1)^2\), Navier BC \(w=\Delta w=0\); beam \((-1,1)\), Navier
BC \(w=w''=0\); modes \((m,m)\), \((m,m{+}1)\), and beam \(m\);
\(\omega_0=\max(10,2\pi f)\), \(\sigma=\max(2,\pi f)\) with
\(f=\max(m,n)\); residual scale \(S^2\).

% ===========================================================================
\subsection{Linearized Korteweg--de Vries / dispersive wave}\label{sec:bm:kdv}

The third-order dispersive operator is exercised by a linearized KdV (Airy)
equation in space--time,
\begin{equation}\label{eq:kdv_pde}
  u_t + \delta\,u_{xxx} = f
  \qquad \text{in } (x,t) \in (-1,1)^2,
  \qquad \delta = 1,
\end{equation}
with the oscillatory travelling profile
\begin{equation}\label{eq:kdv_sol}
  u_\star(x,t) = \sin(k\pi x)\,\cos(k\pi t)
\end{equation}
and Dirichlet data \(u = u_\star\) on the space--time boundary. The source is
\begin{align}\label{eq:kdv_source}
  f
  &= u_{\star,t} + \delta\,u_{\star,xxx} \notag\\
  &= -k\pi\,\sin(k\pi x)\sin(k\pi t)
     - (k\pi)^3\,\cos(k\pi x)\cos(k\pi t),
\end{align}
and the wavenumber is swept over \(k \in \{2,3,4,5,6\}\).

\paragraph{Provenance and comparability.}
KdV is a canonical PINN benchmark; Raissi et al.~\cite{Raissi2019} use the
nonlinear form \(u_t + \lambda_1 u u_x + \lambda_2 u_{xxx} = 0\) with
\(\lambda_2 = 0.0025\). We use the linear dispersive (Airy) operator with a
manufactured oscillatory solution in order to isolate the odd third-order term
\(u_{xxx}\) from the convective nonlinearity; the canonical nonlinear KdV can be
added as a separate case for a direct comparison.

\paragraph{Setting.}
\((x,t)\in(-1,1)^2\); Dirichlet \(u=u_\star\); \(\delta=1\); sweep
\(k\in\{2,3,4,5,6\}\); \(\omega_0=\max(10,2\pi k)\), \(\sigma=\max(2,\pi k)\).

% ===========================================================================
\subsection{Cahn--Hilliard}\label{sec:bm:ch}

The phase-field Cahn--Hilliard equation combines a high-order linear operator
with a cubic nonlinearity. In one space dimension plus time,
\begin{align}\label{eq:ch_pde}
  u_t &= M\big[\Delta(u^3) - \Delta u - \gamma\,\Delta^2 u\big] && (\text{4th order}),\\
  u_t &= M\big[\Delta(u^3) - \Delta u - \gamma\,\Delta^2 u + \kappa\,\Delta^3 u\big] && (\text{6th order}),
  \label{eq:ch_pde6}
\end{align}
with \(M = \gamma = \kappa = 1\), \(\Delta = \partial_x^2\), and
\((x,t) \in (-1,1)^2\). The manufactured solution is
\begin{equation}\label{eq:ch_sol}
  u_\star(x,t) = \sin(a\pi x)\,\cos(a\pi t),
\end{equation}
the nonlinear flux is expanded as in \eqref{eq:lap_u3}, and the source term and
the boundary targets (on \(u\), \(u_{xx}\), and additionally \(u_{xxxx}\) for the
sixth-order case) are obtained from the analytic solution. We use
\(a \in \{2,3\}\) at orders \(4\) and \(6\).

\paragraph{Provenance and comparability.}
Cahn--Hilliard is a standard fourth-order phase-field PINN benchmark, and a
sixth-order variant is used in randomized-smoothing PINNs; chaotic/high-order
space--time PDEs of this type (Kuramoto--Sivashinsky, Gray--Scott) form the
nonlinear high-order tier of PINNacle~\cite{Hao2024PINNacle}. The manufactured
solution \eqref{eq:ch_sol} gives an exactly known reference while retaining the
cubic nonlinearity and the genuine fourth/sixth-order operator.

\paragraph{Setting.}
\((x,t)\in(-1,1)^2\); \(M=\gamma=\kappa=1\); orders \(\{4,6\}\); sweep
\(a\in\{2,3\}\); \(\omega_0=\max(10,2\pi a)\), \(\sigma=\max(2,\pi a)\).

% ===========================================================================
\subsection{Cubic nonlinear Schr\"odinger}\label{sec:bm:nls}

The first genuinely complex-valued benchmark is the cubic nonlinear
Schr\"odinger (NLS) equation,
\begin{equation}\label{eq:nls_pde}
  \mathrm{i}\,u_t + \tfrac12\,u_{xx} + |u|^2 u = f,
  \qquad u:\ \Omega \to \mathbb{C}.
\end{equation}
We use the standing bright-soliton family as the manufactured solution,
\begin{equation}\label{eq:nls_sol}
  u_\star(x,t) = \operatorname{sech}(x)\,\mathrm{e}^{\mathrm{i}kt},
\end{equation}
for which a direct computation gives the source
\begin{align}\label{eq:nls_source}
  \mathrm{i}\,u_{\star,t} + \tfrac12\,u_{\star,xx} + |u_\star|^2 u_\star
  &= \Big(\!-k + \tfrac12 - \operatorname{sech}^2 x + \operatorname{sech}^2 x\Big) u_\star
   = \big(\tfrac12 - k\big)\,u_\star ,
\end{align}
so \(f = (\tfrac12 - k)\,u_\star\), and at \(k = \tfrac12\) the source vanishes
and \(u_\star\) is the exact bright soliton. The temporal frequency is swept over
\(k \in \{1,2,4\}\). A single complex-sinh network carries \(u\) in
\(\mathbb{C}\) natively; the real baselines carry \(u\) as a split real/imaginary
pair, parameter-matched as a pair.

\paragraph{Provenance and comparability.}
This is the complex-valued benchmark introduced by Raissi et
al.~\cite{Raissi2019}, whose canonical setting is
\begin{equation}\label{eq:nls_raissi}
  \mathrm{i}\,h_t + \tfrac12 h_{xx} + |h|^2 h = 0,
  \quad x\in[-5,5],\ t\in[0,\pi/2],\quad
  h(0,x) = 2\operatorname{sech}(x),
\end{equation}
with periodic boundary conditions. We use the same domain window
\(x\in[-5,5]\), \(t\in[0,\pi/2]\), with the manufactured standing soliton
\eqref{eq:nls_sol} in place of the free amplitude-\(2\) evolution, so the
reference solution is known in closed form while the equation, the cubic
nonlinearity, and the complex-valued field are identical to the canonical
setting.

\paragraph{Setting.}
\(x\in[-5,5]\), \(t\in[0,\pi/2]\) (inputs normalized to \([-1,1]^2\); chain-rule
factors \(1/L_x\), \(1/L_t\)); Dirichlet/initial penalty; sweep \(k\in\{1,2,4\}\);
complex vs.\ split-real representation.

% ===========================================================================
\subsection{Time-harmonic Maxwell (lossy medium)}\label{sec:bm:maxwell}

The second complex-valued benchmark is the time-harmonic transverse-magnetic
(TM\(_z\)) Maxwell system, which for the out-of-plane field reduces to a complex
Helmholtz equation,
\begin{equation}\label{eq:maxwell_pde}
  \Delta E + \kappa^2 E = f
  \qquad \text{in } \Omega = (-1,1)^2,
  \qquad \kappa^2 = (a\pi)^2\,(1 + \mathrm{i}\beta),
\end{equation}
with loss tangent \(\beta = 0.2\). The complex permittivity makes the field
genuinely complex-valued, and the manufactured solution is a traveling plane
wave,
\begin{equation}\label{eq:maxwell_sol}
  E_\star(x,y) = \mathrm{e}^{\mathrm{i}a\pi(x+y)},
\end{equation}
with Dirichlet data \(E = E_\star\) on \(\partial\Omega\) and source
\begin{equation}\label{eq:maxwell_source}
  f = \big(-2(a\pi)^2 + \kappa^2\big)\,E_\star .
\end{equation}
The wavenumber is swept over \(a \in \{2,4,6\}\). As in the NLS case, the complex
network carries \(E\) natively while the real baselines use a split-real pair.

\paragraph{Provenance and comparability.}
Time-harmonic Maxwell / frequency-domain electromagnetic scattering is an
inherently complex-valued PDE family for which complex-valued surrogates have
shown clear advantages over real-valued ones~\cite{Jiang2024Maxwell}. The lossy
(\(\mathrm{i}\beta\)) term is the linear complex-valued companion to the cubic
NLS test: it forces a genuinely complex field through a linear second-order
operator.

\paragraph{Setting.}
\(\Omega=(-1,1)^2\); Dirichlet \(E=E_\star\); loss tangent \(\beta=0.2\); sweep
\(a\in\{2,4,6\}\); \(\omega_0=\max(10,2\pi a)\), \(\sigma=\max(2,\pi a)\);
complex vs.\ split-real representation.
